\SourceLeaf{319}%
up to but excluding fifth order, with \(s\) treated as a quantity of first order, and that in the parenthesized terms on the right, \(q,\alpha l\) and \(u\) may be replaced by \(s\mathord{.}\cos U\), \(s\mathord{.}\frac{\sin U}{\cos Q}\), \(s\mathord{.}\sin U\mathord{.}\operatorname{tang}Q\) without loss of accuracy.

\Article{22}
The quantities \(B,b,T,t,r\), which refer to the ellipsoid, must first be compared with their counterparts on the sphere, \(Q,q,U,u,As\). All the equations to be set up for this purpose will be accurate through at least the third order inclusive. The derivation itself will make it easy to see that this requirement is met.

If the series given in Art. 8 is applied to our two points, the quantities denoted there in general by \(p\) and \(q\) must be expressed in our present notation
\begin{center}
\begin{tabular}{ll}
for the first point by & \(B+\frac12b-P\) and \(\frac12q\)\\
for the second point by & \(B-\frac12b-P\) and \(-\frac12q\)
\end{tabular}
\end{center}
and we therefore have the two equations
\begingroup\fontsize{10}{12}\selectfont
\[
\begin{aligned}
B+\tfrac12b={}&P+\frac{\cos\theta}{2\cos\varphi}\mathord{.}q-\frac{3ee}{8\cos\varphi^2}\mathord{.}\cos P\mathord{.}\sin P\mathord{.}qq\\
&+\frac{ee}{16\cos\varphi^3\cos\theta}(-\cos P^2+\sin P^2+ee(5\cos P^2\mathord{.}\sin P^2-\sin P^4))q^3\\
B-\tfrac12b={}&P-\frac{\cos\theta}{2\cos\varphi}\mathord{.}q-\frac{3ee}{8\cos\varphi^2}\mathord{.}\cos P\mathord{.}\sin P\mathord{.}qq\\
&-\frac{ee}{16\cos\varphi^3\cos\theta}(-\cos P^2+\sin P^2+ee(5\cos P^2\mathord{.}\sin P^2-\sin P^4))q^3
\end{aligned}
\]
\endgroup
Addition and subtraction thus give
\[
\begin{aligned}
B={}&P-\frac{3ee}{8\cos\varphi^2}\mathord{.}\cos P\mathord{.}\sin P\mathord{.}qq\quad .\ .\ .\ .\ .\ .\ .\ .\ .\ .\ .\ .\ .\ .\ .\ .\ .\ .\ .\quad(4)\\
b={}&\frac{\cos\theta}{\cos\varphi}\mathord{.}q+\frac{ee}{8\cos\varphi^3\cos\theta}(-\cos P^2+\sin P^2\\
&\qquad+ee(5\cos P^2\mathord{.}\sin P^2-\sin P^4))q^3\qquad(5)
\end{aligned}
\]
It is easy to see, moreover, that equation (4) differs from the exact relation by fourth-order quantities, while equation (5) differs only by fifth-order quantities.

To compare \(T\) and \(t\) with \(U\) and \(u\), the formulas developed at the end of Art. 15 must be used. They were based on an assumption
